% ==================== 第 1 章 作品概述 ====================
\section{作品概述}

\subsection{作品基本信息}

\noindent 注：正式提交时必须贴附盖章版本的报名表第一页和第二页截图。

\begin{longtable}{@{}p{5.0cm}p{10.0cm}@{}}
\toprule
信息项 & 填写内容 \\
\midrule
挑战杯系统报名时参赛作品名称（必填） & 基于千问大模型的多角色协同 AI Scientist 系统 \\
最终参赛作品名称（如有更新可备注） & AstroQuery AI：基于 Qwen 多模态大模型与可解释闭环的天文科学数据智能挖掘与跨源整合系统 \\
参赛选题（必填） & 赛道二-选题 1-A \\
参赛作品简介（300 字内，必填） & 面向天文科研数据获取的智能检索系统：研究者输入自然语言查询目标天体参数（如“昴星团的年龄、距离和金属丰度”），系统自动完成意图解析与目标确认，并行检索多源专业星表与海量文献库。依托多模态大模型，系统从论文正文、复杂表格及科学图表中提取相关数值与图文佐证，并为每条数据生成可精确定位回溯的原文证据链。随后数据进入闭环质检流水线：多维校验单位格式与物理合理性，给出置信评分；自动规范清洗并全程修改留痕；对多源数据展开冲突分析与分歧归因，疑难异常流转人工复核。最终结合字段洞察在前端可视化呈现，并支持导出高可信度的结构化科研数据。 \\
参赛作品使用 Qwen 大模型及 AI 技术说明（300 字内，必填） & 本作品依托阿里云百炼平台调用通义千问系列大模型，构建快模型分流、强模型把关与多模态模型提取的三级协同体系。在全流程中，由高精度的 qwen3.8-max 结合天文知识筛选物理性质并确立全局字段标准，由 qwen3.7-plus 多模态模型逐页深度解析论文正文与图表以精准提取数值、单位和原文依据，由高响应的 qwen3.7-flash 贯穿负责意图澄清、检索式生成、定位兜底、质量评估与冲突分类。系统通过白名单约束防范幻觉，坚持模型判断与规则执行分离，杜绝数值错误改写，保障天文科学数据的严谨可靠与全程可溯源。 \\
宣传视频链接/其他作品相关文件链接（选填） & \emph{（视频链接待补充：夸克网盘分享链接）} \\
\bottomrule
\end{longtable}

\begin{figure}[H]\centering
\includegraphics[width=0.47\textwidth]{\detokenize{报名表1.png}}\hfill
\includegraphics[width=0.47\textwidth]{\detokenize{报名表2.png}}
\caption{盖章报名表第一页（左）与第二页（右）}\label{fig:signup}
\end{figure}

\subsection{本作品核心主张}

\begin{itemize}[leftmargin=1.2em,itemsep=2pt]
\item \textbf{本作品服务的具体科研数据需求：}天文研究在深入剖析特定恒星、星团或星系等天体时，亟需一次性全面获取单一目标在多个物理性质上的海量数据。针对单个天体，科研人员往往需要跨越海量文献与星表，深度收集多种性质的数十乃至上百条分散记录。受限于天体别名繁杂、图表分散、单位不一以及测量差异，人工逐篇筛查耗时极长且极易遗漏。科研工作迫切需要针对目标天体实现多性质的深度挖掘与规范对齐，快速产出高密度、多维度、可溯源的结构化数据。
\item \textbf{本作品实际完成的核心方法：}系统通过人机交互确认目标天体与所需性质，结合天文数据库展开全别名，并依据知识库生成统一性质清单；随后并行检索专业星表与文献论文，利用多模态大模型解析正文与图表提取数值，配合文本层逆向定位提供精准原文位置；最后通过质量流水线对字段格式与合理范围进行多维质检，由确定性规则完成误差区间拆解与单位规范化，并利用离群统计分析自动归因多源测量冲突，在保留原始分歧与全流程修改留痕的前提下完成结构化交付。
\item \textbf{最有代表性的结果 1（数据或任务范围、采用的方法、实际结果）：}疏散星团昴星团 M45 多维物理参数智能检索与跨源整合。
\begin{itemize}[leftmargin=1.2em,itemsep=1pt]
\item \textbf{数据与任务范围}：针对疏散星团昴星团 M45，输入年龄、距离与金属丰度查询词条。
\item \textbf{采用的方法}：采用需求意图澄清、性质白名单构建、星表文献并行检索、多模态图文提取与文本层逆向定位，结合质量管线进行确定性规则清洗与离群冲突分类的全流程自动化链路。
\item \textbf{实际结果}：系统自动将查询扩展细化为日心距离、视差、距离模数、年龄与金属丰度 5 项规范化物理量；跨越 1998 至 2025 年间 10 种天文期刊，自动检索并融合 29 篇论文与 MWSC 星表，抽取整合 103 条多源记录；五项物理量分别获得 37 条距离记录集中在 134.4 至 136.2 pc，25 条距离模数记录集中在 5.58 至 5.65 mag，7 条视差记录集中在 7.42 至 7.49 mas，26 条年龄记录集中在 100 至 130 Myr，以及 8 条金属丰度记录集中在 $-$0.03 至 +0.03 dex；生成 78 张图表证据，正文定位偏差中位数为 0.02 倍字宽，表格框选覆盖率达百分之百；系统对早期依巴谷测距偏差与对数单位写法等 31 项跨源冲突进行了分类标注并保留原始记录；经 69 条修改留痕，综合质量得分达 0.85，全程耗时约 21 分钟自动完成。
\end{itemize}
\item \textbf{最有代表性的结果 2（数据或任务范围、采用的方法、实际结果）：}单一恒星“北河三（Pollux）”多性质检索与可解释的跨源一致性整合。
\begin{itemize}[leftmargin=1.2em,itemsep=1pt]
\item \textbf{数据与任务范围}：针对点源恒星北河三（Pollux），输入质量、半径、表面温度与光谱型查询词条。
\item \textbf{采用的方法}：沿用全流程自动化链路，针对单星目标适配多模态检索、物理规律跨字段自洽校验与显式可解释归因分析。
\item \textbf{实际结果}：系统跨越 2006 至 2026 年间 ApJ、A\&A、MNRAS 等 3 种核心期刊（7 篇文献）以及 HD、Hipparcos、HR、TIC 4 个权威星表，自动检索并融合 11 个数据源，抽取整合 16 条多源记录；四项物理量分别获得 5 条质量记录（1.7 至 2.31 Msun，文献共识收敛于 1.9 至 2.1 Msun）、3 条半径记录（8.8 至 9.3 Rsun 并经双方法互证）、4 条表面温度记录（4892.6 至 5000 K）以及 4 条光谱型记录（统一规范归入 K0 III 巨星族）；系统对初始/现质量口径差异及测光/谱综合标定系统差进行了显式归因解读，利用 Stefan-Boltzmann 定律与 HR 图完成多参数跨字段自洽校验，并对星表缺失字段保留显式标记；经 6 处规范化修改留痕，实现零未解冲突与免人工复核，综合质量得分达 0.9235（Excellent），全程耗时约 14 分钟自动完成。
\end{itemize}
\item \textbf{目前仍存在的主要局限：}① \textbf{受学术出版版权壁垒客观制约}，目前全文本抽取主要依赖 arXiv 预印本及 Open Access（Unpaywall）渠道，封闭期刊的直接解析覆盖面仍受客观生态限制；② \textbf{性质库覆盖面有边界}——当前覆盖主流星表及高频物理量，对少数前沿、长尾或非标准化细分性质的自适应推断与泛化映射能力仍需进一步增强；③ \textbf{评估算法难以完全客观}——质量评分、冲突判定等自动化评估依赖预设规则与模型判断，其客观性受规则完备性与模型能力制约，极端情况下仍需人工复核兜底。
\end{itemize}

\subsection{团队实际解决的问题}

在天体物理研究中，获取目标天体关键物理参数是开展建模与理论验证的基础，但现有获取方式存在多重不足：

\renewcommand{\labelenumi}{\arabic{enumi}．}
\begin{enumerate}[leftmargin=1.6em,itemsep=2pt]
\item \textbf{文献暗数据滞后}：大量最新拟合参数作为暗数据深嵌于论文正文与复杂表格中，数据库收录往往滞后数月乃至一年以上；
\item \textbf{别名漏检}：天体跨巡天跨文献存在数十种别名，人工逐一核对检索极易漏检关键成果；
\item \textbf{单位与量纲异构}：物理量表述形式与量纲基准异构多样，迫使科研人员消耗大量精力在单位换算与手工对齐上；
\item \textbf{摘录丢出处、通用模型不可信}：人工摘录易丢失出处依据，而直接使用通用大模型又面临不可接受的数值幻觉且缺乏原文位置佐证。
\end{enumerate}

\subsection{本作品实际完成的工作}

本作品已实现从科研需求到数据交付的六阶段全流程闭环功能：

\begin{enumerate}[leftmargin=1.6em,itemsep=2pt]
\item \textbf{需求解析}：通过自然语言意图理解与人机多轮澄清明确目标，结合天文数据库统一标识并展开全别名，依据领域知识库生成统一性质清单；
\item \textbf{检索}：基于性质清单并行检索多张专业星表并自动生成文献检索式，实现论文及附表数据的靶向筛选与下载；
\item \textbf{提取}：依托多模态大模型逐页解析论文正文与复杂表格以提取物理量数值，配合文本层逆向定位锁定真实坐标并筛选图表证据；
\item \textbf{质检}：通过多维评估打分、确定性规则清洗与离群冲突检测构建质控流水线，完成误差区间拆解与单位规范化，对多源测量冲突进行显式归因并保留原始分歧，支持状态机闭环迭代与人机协同裁决；
\item \textbf{数据产出}：输出标准化数据集并完整保留原始出处、页码坐标、逐字摘录与修改留痕，实现全要素深度可溯源；
\item \textbf{前端显示}：提供多维数据检索、原文高亮定位回溯、图证画廊展示、冲突洞察分析与任务流交互回放。
\end{enumerate}

\subsection{本作品形成的完整闭环与主要成果}

\begin{figure}[H]\centering
\includegraphics[width=0.96\textwidth]{\detokenize{real_process.jpg}}
\caption{真实查询全流程展示——从科研数据需求到结构化数据的总体思路}\label{fig:reallflow}
\end{figure}

\begin{itemize}[leftmargin=1.2em,itemsep=2pt]
\item \textbf{本作品最终形成的主要数据成果：}每次查询产出带来源追溯的结构化数据包，包括结构化数据文件、长表与宽表、质量评分、来源信息、修改记录与洞察报告，以及论文图表证据和原文定位页面，供研究者下载与核验。
\item \textbf{本作品已经实际处理的数据来源和任务范围：}数据来源涵盖专业星表、文献检索库、开放获取论文及其附带数据表；任务范围覆盖恒星个体、星团与星协、河外星系与活动星系核、星系群与星系团 4 大天体类别约 100 种细分天体，构建了每类包含 20--50 种性质的参数体系适配查询，统一支持星等、距离、温度、光度、质量、半径、金属丰度等跨类共用的核心物理量。
\item \textbf{相较通用检索或人工复制整理，本作品最主要的不同：}通用检索只返回文档列表，人工整理费时且难以追溯；本作品打通需求理解、查找、解析、对齐、整合、质检、输出的完整链路，以统一的性质清单约束所有提取路径，每条记录的精确出处保存在来源信息里，多源冲突经系统化判断与标注而非简单合并，最终输出可复现、可追溯、可下载的结构化数据。
\end{itemize}

% ==================== 第 2 章 数据需求与来源 ====================
\section{数据需求与来源}

\subsection{开场：一次真实查询}

做疏散星团研究的科研人员在校准恒星演化模型时，常常要为一个目标天体凑齐“年龄、距离、金属丰度”这类基础参数——而它们散落在几十篇文献与若干星表里。他输入了一句再平常不过的话：\textbf{“昴星团的年龄、距离和金属丰度”}。本作品的一次完整查询就此开始。

本章回答两件事：\textbf{这单查询背后“需要哪些数据”是如何被确定的}（需求规格化），以及\textbf{哪些来源能满足这些需求}（来源体系）。昴星团（M45）作为全报告贯穿主线——第 3 章各模块的实况卡、第 4 章的版本迭代与第 6 章的代表档案都以这次真实查询为例——它命中数据库、论文正文、论文表格、论文图证四类来源，覆盖 5 项物理性质与 1998--2025 年 28 年文献跨度，并天然携带一段真实的测距争议史，足以把系统的能力与边界讲透。

\subsubsection{为何以昴星团 M45 作为贯穿展示对象}

\begin{itemize}[leftmargin=1.2em,itemsep=2pt]
\item \textbf{案例对应的研究问题或数据使用场景：}疏散星团基准参数测定——昴星团（M45）是天文学中经典的基准星团。研究人员在校准恒星演化理论模型、测量银河系距离尺度，以及研究年轻恒星与行星系统演化等实际科研工作中，都需要获取其精确的年龄、距离和金属丰度，作为计算与分析的基础输入数据。
\item \textbf{选择它进行完整展示的实际理由：}该案例一次命中多种来源类型（数据库、论文正文、论文表格、论文图证），性质覆盖广（5 项），文献跨度 28 年（1998--2025），且存在真实的数据争议（早期与近代测距结果分歧），能够完整呈现从查找、解析到质量校验的全流程。
\item \textbf{它能够覆盖哪些来源类型和关键能力：}覆盖数据库（星团目录）、论文正文、论文表格、论文图证四类来源的实际解析；覆盖多轮澄清、性质清单约束、数值与单位分离提取、原文定位、单位统一、冲突保留与洞察报告等关键能力。
\item \textbf{该案例不能代表哪些任务或条件（边界声明）：}不能代表单一恒星（点源）任务——此类任务另见 1.2 结果 2（Pollux）；不能代表时变性质（如光变周期）、论文补充材料命中（本案例未命中补充材料源）以及封闭期刊无开放全文的场景。
\end{itemize}

\subsection{需求规格化：这项研究实际需要什么数据}

模板设问面向“一项具体研究”；本作品是面向任意自然语言查询的工具，因此以贯穿案例 M45 实填五问，展示一次查询如何被规格化为一套明确的数据需求。

\begin{longtable}{@{}p{4.4cm}p{10.6cm}@{}}
\toprule
要回答的问题 & 团队实际填写内容（M45 实例） \\
\midrule
研究要解决什么问题，或者这些数据最终用于哪项分析 & 为昴星团（M45）的星团物理研究提供多源一致的结构化参数：年龄、距离、距离模数、视差、金属丰度用于恒星演化模型校验与疏散星团距离尺度研究（该星团是恒星演化与测距的关键标定对象，文献报道量极大且存在距离争议）。 \\
需要收集哪些对象、材料、记录或信息 & 目标天体 M45 的宏观物理参数，来源包括：① 数据库与星表记录（星团目录的年龄、距离、金属丰度）；② 论文正文与表格中报告的值（29 篇独立论文，含从文本与表格提取的数值）；③ 论文中的图表证据（78 张，作为证据链）。 \\
需要收集哪些范围内的数据（如样本范围、实验条件、时间范围或研究区域） & M45 单一天体；字段范围限定为年龄（yr）、距离（pc）、距离模数（mag）、视差（mas）、金属丰度（dex）5 个性质；文献时间跨度为各来源论文的发表年份（1998--2025）。 \\
每条数据需要包含哪些内容或字段；如涉及数值，再写单位和精度 & 每条记录包含：物理量名称、数值（如 135.15、125 $\pm$ 20）、单位（如 yr、pc、mag、dex 等）、来源、提取方式（星表查询或论文提取）、论文页码与原文位置、置信度、原始写法，用于溯源核对。 \\
最终要交付什么格式的数据，后续如何使用 & 结构化数据文件与长表宽表，附质量报告、来源信息与洞察报告；可直接用于参数汇总、多论文对比或作为后续恒星演化分析的输入。 \\
\bottomrule
\end{longtable}

\noindent\textbf{性质清单（字段白名单示例，本案例实际生成的 5 项）}：每条记录在提取前受此清单约束，字段名、标准单位与含义全局唯一。

\begin{longtable}{@{}p{2.4cm}p{2.6cm}p{2.2cm}p{6.8cm}@{}}
\toprule
标准字段 & 含义 & 标准单位 & 本案例各来源写法示例（原文保留） \\
\midrule
age & 星团年龄 & yr & Myr、yr、$10^{6}$ yr、log(yr) \\
distance & 日心距离 & pc & pc \\
dist\_modulus & 距离模数 & mag & mag \\
parallax & 视差 & mas & mas \\
fe\_h & 金属丰度 & dex & dex、log(Sun) \\
\bottomrule
\end{longtable}

\noindent 该清单即提取环节的白名单（见 3.3.2），也是质量评价判据之一（见 4.1 与 3.4.1），全文同源。

\medskip
\noindent\textbf{推广段（任意查询如何规格化）}：任意自然语言查询都经历同款规格化流程：目标天体由天文数据库解析为标准名称与别名，性质由领域知识库与强模型按天体类型筛选。系统任务范围覆盖恒星个体、星团与星协、河外星系与活动星系核、星系群与星系团 4 大天体类别约 100 种细分天体，每类预置 20--50 种性质的参数体系，并统一支持星等、距离、温度、光度、质量、半径、金属丰度等跨类共用的核心物理量（详细流程见 3.1 架构与 3.3 检索提取）。

\subsection{来源体系：实际使用的数据与资料来源}

上述数据需求由四类来源互补满足：\textbf{天体标识数据库}解决“目标是谁、还有哪些名字”的身份问题，\textbf{专业星表}提供字段规范、覆盖大样本的编目测量值，\textbf{学术文献}承载数据库收录滞后数月乃至数年的前沿数值（含图表证据），\textbf{文献附带数据表}补上正文未逐条列出的完整测量序列。需要说明的是，四类来源在“同一天体的同一性质”上必然存在覆盖、重复与写法/取值差异——这并非缺陷，而是交叉印证与争议呈现的原料；其系统化的去重、统一与冲突裁决在第 3.4 节质量管线中处理，此处先给出来源总览与其关系保留承诺。

\begin{longtable}{@{}p{2.7cm}p{3.2cm}p{3.4cm}p{3.2cm}p{2.5cm}@{}}
\caption{实际使用的四类数据来源}\label{tab:sources}\\
\toprule
数据或资料来源 & 本作品实际获取的内容 & 选择该来源的理由 & 覆盖范围与已知局限 & 来源关系如何保留 \\
\midrule
\endfirsthead
\toprule
数据或资料来源 & 本作品实际获取的内容 & 选择该来源的理由 & 覆盖范围与已知局限 & 来源关系如何保留 \\
\midrule
\endhead
\textbf{天体标识数据库}（如 SIMBAD） & 目标天体的标准中文/英文名称、天体类型分类、天球坐标及全量别名列表 & 作为天体身份的统一锚点，利用全别名清单解决跨文献检索时的漏检问题 & 覆盖绝大多数已知天体；中文名称需先转写为国际标准标识 & 记录为全局天体元数据，作为唯一样本标识贯穿检索、提取与输出全链路 \\
\textbf{专业天文星表}（25 个，含各类星团星表与综合目录） & 视差、自行、多波段星等、距离模数、年龄、金属丰度等目录级编目测量值 & 属于权威编目数据，字段格式规范，覆盖大样本多巡天观测历史 & 仅提供目录级汇总参数，无文献讨论细节；不同星表的字段命名方式存在差异 & 记录所在数据表名称、检索键列名、行号及原始单位，支持精准反查 \\
\textbf{学术文献}（论文正文、表格及内嵌科学图表） & 正文、表格中的最新观测或拟合数值，与研究内容相关的图表图片 & 论文是前沿物理数值的主体来源，包含数据库尚未收录的最新研究成果；图表提供直观证据 & 受版权限制目前主要支持开放获取（OA）及预印本文献，封闭期刊正文无法直接抓取 & 记录文献唯一检索编号（DOI / 文献标识码）、发表年份、作者、PDF 页码、正文与表格定位框坐标及关联图表切片 \\
\textbf{文献附带数据表}（期刊补充材料 / CDS 附表） & 样本恒星测光表、详细参数拟合表等高密度结构化数据 & 数据表通常包含正文未逐一列出的完整原始测量序列，是对正文的重要补充 & 复杂排版与跨页数据表存在识别风险，无法确定列属性时整表丢弃 & 记录数据表编号、所在文件页码、行号索引以及表头原始列名 \\
\bottomrule
\end{longtable}

\noindent 来源间覆盖/重复/冲突的细节不在此展开，其系统化处理见 3.4.2 与 3.4.3；“如何判断某一来源可用”“无法核实的数据如何处理”两问属检索筛选规则，已移入 3.3.1 检索节。

\subsection{系统承诺边界：能做到和不能做到的范围}

\begin{itemize}[leftmargin=1.2em,itemsep=2pt]
\item \textbf{能够查找和解析哪些实际资料（论文、数据库、网页、表格、附件或图表等）：}论文正文与复杂表格、论文图表、专业星表、文献库元数据、论文附带的数据表。
\item \textbf{输出中哪些数据直接来自原始资料，哪些经过规则换算或模型补充：}① \textbf{直接来自原始资料}：从论文原文逐字提取的数值（附原文摘录佐证）与星表查询的原始记录；② \textbf{规则换算}：单位统一换算，全部按既定规则执行并留下修改记录；③ \textbf{模型补充}：数值格式清理、跨字段关系分析与使用建议等，均明确标注为模型生成。
\item \textbf{找不到数据、不同来源冲突或无法确认时，系统实际怎样呈现：}找不到数据时，前端记录表格中该字段不出现记录，质量报告与洞察报告会明确列出缺失字段（如部分星表缺少质量、半径、温度字段会被标记为缺失项）；不同来源数值冲突时，前端记录详情中展示冲突标注，写明涉及来源、取值区间与处置方式，质量卡片同步显示冲突数量；无法确认是否为同一对象的数据以待审核状态呈现，等待研究者决策，系统不做自动合并。
\item \textbf{哪些内容仍需要研究者人工核对：}跨来源写法差异（如同一量在不同文献中以不同单位呈现）是系统的刻意设计——原文写法全部保留并标注，最终取舍交由研究者按领域知识决定；普通异常值在洞察报告中给出成因假设供研究者参考，无需人工介入。系统仅在自动判断确属无法裁决时发起人机协同复核：某来源数据整体不可信（提取质量极差、超合理范围的记录过半）、跨来源实体错配无法自动区分、或处理环节异常兜底——此时前端向研究者发起确认，研究者做出选择后流程继续，系统不擅自合并或丢弃。
\end{itemize}
